\section{Lý do chọn đề tài}
    Trong các hệ thống Internet of Things (IoT) hiện nay, IoT Gateway giữ vai trò trung tâm trong việc kết nối các thiết bị đầu cuối với hạ tầng mạng và nền tảng đám mây. Gateway không chỉ thực hiện chức năng chuyển tiếp dữ liệu mà còn đảm nhiệm các tác vụ quan trọng như chuyển đổi giao thức, xử lý dữ liệu tại biên, quản lý thiết bị và thực thi các cơ chế bảo mật.

    Tuy nhiên, thực tế triển khai cho thấy mỗi hệ thống IoT lại có yêu cầu khác nhau về loại thiết bị, chuẩn giao tiếp và giao thức truyền thông. Việc thiết kế hoặc lựa chọn một gateway cố định, tích hợp sẵn nhiều giao thức không cần thiết dẫn đến lãng phí tài nguyên, tăng chi phí phần cứng và giảm tính linh hoạt khi mở rộng hoặc thay đổi thành phần của hệ thống.

    Mặt khác, các gateway thương mại hiện nay thường đi theo hai xu hướng: hoặc tích hợp quá nhiều chức năng dẫn đến giá thành cao và phức tạp trong vận hành, hoặc quá đơn giản, khó mở rộng và không đáp ứng được các bài toán IoT đa dạng trong thực tế. Ngoài ra, việc cấu hình gateway trong nhiều trường hợp vẫn phụ thuộc nhiều vào việc chỉnh sửa firmware, gây khó khăn cho quá trình triển khai và bảo trì tại hiện trường.

    Xuất phát từ những vấn đề trên, nhóm lựa chọn đề tài “Thiết kế và hiện thực một gateway IoT dạng mô-đun có khả năng cấu hình” với định hướng xây dựng một giải pháp kỹ thuật linh hoạt, có thể tùy biến theo nhu cầu sử dụng, tối ưu chi phí và phù hợp với các kịch bản triển khai IoT thực tế.